\begin{trapbox}
	\begin{description}[style=nextline, leftmargin=2.8em, labelsep=0.5em, itemsep=0.35em]
		\item[\textbf{\textcolor{CTrap}{易错点一（平移量混淆）}}] 先伸缩后平移时，学生常将平移量误算为 $|\varphi|$，导致得到错误解析式。
		\item[\textbf{\textcolor{CTrap}{正确理解（正确算法）：}}] 应先伸缩后平移，平移量应为 $|\varphi/\omega|$，即 $\left|\frac{\varphi}{\omega}\right|$。
		\item[\textbf{\textcolor{CTrap}{错因分析（根源问题）：}}] 错误在于忽略了横向伸缩已经改变了横坐标；反过来，先平移后伸缩也不能误用 $|\varphi/\omega|$，该路径应平移 $|\varphi|$。
	\end{description}

	\medskip
	\begin{description}[style=nextline, leftmargin=2.8em, labelsep=0.5em, itemsep=0.35em]
		\item[\textbf{\textcolor{CTrap}{易错点二（方向混淆）}}] 学生认为 $\varphi>0$ 向右平移，$\varphi<0$ 向左平移，实际相反。
		\item[\textbf{\textcolor{CTrap}{正确理解（左加右减）：}}] $\varphi>0$ 时向左平移，$\varphi<0$ 时向右平移，$\varphi=0$ 时不平移。两条路径距离分别为 $|\varphi|$ 和 $|\varphi|/\omega$，方向相同。
		\item[\textbf{\textcolor{CTrap}{错因分析（思维误区）：}}] 将 $x+\varphi$ 直观理解为“加就是向右”，混淆了自变量变换与点坐标移动的关系。
	\end{description}

	\medskip
	\begin{description}[style=nextline, leftmargin=2.8em, labelsep=0.5em, itemsep=0.35em]
		\item[\textbf{\textcolor{CTrap}{易错点三（顺序僵化）}}] 部分学生只熟悉一种变换顺序，遇到另一种时不会调整平移量。
		\item[\textbf{\textcolor{CTrap}{正确理解（灵活运用）：}}] 两条路径在调整平移量后才得到相同终点：先移后缩用有符号距离 $\varphi$；先缩后移用 $\varphi/\omega$。当 $\varphi=0$ 或 $\omega=1$ 时，距离恰好相等。
		\item[\textbf{\textcolor{CTrap}{错因分析（缺乏本质理解）：}}] 误以为水平平移与伸缩一般可交换。实际上 $(x-\varphi)/\omega=x/\omega-\varphi/\omega$，不调整第二条路径的平移量就不能一般地交换顺序。
	\end{description}

	\medskip
	\begin{description}[style=nextline, leftmargin=2.8em, labelsep=0.5em, itemsep=0.35em]
		\item[\textbf{\textcolor{CTrap}{易错点四（负系数）：}}] 把带符号的 $A$ 一律叫作振幅，或只把纵坐标乘以 $|A|$ 就认为得到目标图像。
		\item[\textbf{\textcolor{CTrap}{正确理解（翻折）：}}] 振幅是 $|A|$；最后纵坐标乘以 $A$。若 $A<0$，可先乘以 $|A|$ 再关于 $x$ 轴对称，例如 $-2\sin\theta$ 在 $\sin\theta=1$ 处取值 $-2$。
		\item[\textbf{\textcolor{CTrap}{易错点五（伸缩方向）：}}] 只因 $\omega>0$ 就把横坐标乘以 $1/\omega$ 一律称为压缩。
		\item[\textbf{\textcolor{CTrap}{正确理解（分情况）：}}] $\omega>1$ 是横向压缩，$\omega=1$ 尺度不变，$0<\omega<1$ 是横向拉伸。
	\end{description}
\end{trapbox}
